\section{Introduction}\label{sec:intro}
\subsection{Background and the stability problem}
The study of elastic curves begins with the Euler--Bernoulli bending energy, whose critical points under a length constraint are the classical elasticae. The geometry of closed elasticae and the variational role of their squared curvature were developed systematically by Langer and Singer \cite{LS84}; their subsequent work on curve straightening connected the variational problem with an evolution toward critical curves \cite{LS85}. These two viewpoints lead naturally to a dynamical question: which geometric profiles attract nearby curves under the gradient flow of bending energy?

For a smooth oriented immersion $\gamma:S^1\to\R^2$, let $\dd s=|\gamma_u|\dd u$ denote arclength measure and $\partial_s=|\gamma_u|^{-1}\partial_u$ the arclength derivative. Write $\tau=\partial_s\gamma$ for the unit tangent and $\nu=\operatorname{rot}\tau$ for the unit normal, where $\operatorname{rot}$ is counterclockwise rotation through $\pi/2$. The curvature vector is $\kappa=\partial_s^2\gamma=k\nu$, and the signed curvature is $k=\ip\kappa\nu$. We take $S^1=\R/(2\pi\Z)$, using a $2\pi$-periodic parameter for \eqref{eq:beta}, and write $L[\gamma]=\int_\gamma\dd s$ for the length. For a vector field $v$ along a curve, $v^\perp=\ip v\nu\nu$ denotes its normal projection. The free elastic flow is
\begin{equation}\label{eq:flow}
 (\partial_t\gamma)^\perp=-F_\gamma\nu,
 \qquad F_\gamma=\partial_{s_\gamma}^2k_\gamma+\tfrac12k_\gamma^3.
\end{equation}
It is the $L^2(\dd s)$ gradient flow of $E[\gamma]=\frac12\int_\gamma k^2\dd s$, considered up to orientation-preserving reparametrization. Dziuk, Kuwert and Sch\"atzle proved global existence for smooth closed initial curves \cite[Theorem~3.2]{DKS}; we use the geometric existence and uniqueness statement recorded in \cite[Theorem~1.1]{MW}. Our geometric notation follows \cite[Section~2.1 and Theorem~1.2]{MW}: the absolute turning number is $\omega(\gamma)=\left|(2\pi)^{-1}\int_\gamma k\dd s\right|$. For a time-dependent curve, write $L(t)=L[\gamma(\cdot,t)]$ and $\eta(\cdot,t)=\gamma(\cdot,t)/L(t)$ for the length-normalized flow.

Length constraints and length penalization give a well-developed convergence theory. In particular, Mantegazza and Pozzetta proved smooth convergence for the flow of $E+L$ by a \L{}ojasiewicz--Simon inequality \cite[Theorem~1.2]{MP}. The free problem has a different asymptotic structure: $E[\rho\gamma]=\rho^{-1}E[\gamma]$, and the known closed self-similar solutions expand. Studying an attracting shape therefore requires a normalization that turns an expander into an equilibrium.

Miura and Wheeler established stability of expanding circles, including multiple covers, under a scale-invariant smallness condition on the arclength derivative of curvature \cite[Theorem~1.2]{MW}. They also identified stability of the turning-number-zero Bernoulli lemniscate as an open problem \cite[Section~1.2]{MW}. More recently, Andrews and Wheeler constructed infinitely many geometrically distinct closed homothetic expanders and conjectured that the dynamically stable free-elastic profiles are precisely multiply-covered circles and Bernoulli lemniscates \cite[Theorem~1.1 and Conjecture~1.4]{Jellyfish}. Their work is cited here as a preprint. The result below establishes the lemniscate stability part of that conjectural picture; classification of all stable expanders remains a separate question.

The reference profile is
\begin{equation}\label{eq:beta}
 \beta_{\rm B}(u)=\frac{(\cos u,\sin u\cos u)}{1+\sin^2u},\qquad u\in S^1,
\end{equation}
It has turning number zero and expands homothetically under \eqref{eq:flow} \cite[Section~2.2.3]{MW}. Its signed curvature changes sign, and its self-intersection must be treated on the parameter circle. Multiple covers introduce a further issue: perturbations can vary from one traversal to the next. A proof restricted to the symmetry of the reference curve would not control these additional modes. Our argument treats the entire spectrum on every fixed covering circle.

\subsection{Geometric formulation and nonlinear stability}
Write $\T_\ell=\R/(\ell\Z)$, for $\ell>0$, and
\begin{equation}\label{eq:lengths}
 K=K(1/2)=\int_0^{\pi/2}(1-\tfrac12\sin^2\theta)^{-1/2}\dd\theta,
 \qquad \ell_0=4K/\sqrt2,\quad \ell_M=M\ell_0.
\end{equation}
The arclength-parametrized $M$-fold cover $\beta_M:\T_{\ell_M}\to\R^2$, starting at $u=0$ and following increasing $u$, has curvature
\begin{equation}\label{eq:reference-curvature}
 k_{\beta_M}(s)=3\cn(\sqrt2s\mid1/2).
\end{equation}
Here $1/2$ is the elliptic parameter. The covering multiplicity $M$ is distinct from the turning number: $\omega(\beta_M)=0$ for every $M$. For $a=L[\gamma]/\ell_M$, parametrize the length-matched curve $\widetilde\gamma=a^{-1}\gamma$ by arclength and set
\begin{equation}\label{eq:profile}
 \KM(\gamma)=\ell_M\min_{\sigma\in\T_{\ell_M}}
 \int_0^{\ell_M}|k_{\widetilde\gamma}(s)-k_{\beta_M}(s+\sigma)|^2\dd s.
\end{equation}
This is invariant under translations, rotations, positive dilations and orientation-preserving reparametrizations. Since $\int_0^{\ell_M}k_{\beta_M}\dd s=0$, Cauchy--Schwarz gives
\begin{equation}\label{eq:turn-lock}
 4\pi^2\omega(\gamma)^2\le\KM(\gamma),
\end{equation}
The integer $\omega(\gamma)$ is the absolute degree of the tangent map $\tau:S^1\to S^1$. Thus $\KM(\gamma)<4\pi^2$ forces $\omega(\gamma)=0$.


The functional \eqref{eq:profile} compares the full signed-curvature profile after matching length and arclength phase. It does not require an a priori bound on derivatives of the curvature deviation. Lemma~\ref{lem:curvature-entry} shows that its square root controls an $H^2$ normal graph after geometric alignment. This is the initial regularity at which we enter the normalized evolution. We write $SO(2)$ for the orientation-preserving orthogonal linear maps of $\R^2$.

\begin{theorem}[Local stability]\label{thm:main}
Fix an integer $M\ge1$. There is $\varepsilon_M>0$ such that, for every smooth oriented closed immersion $\gamma_0$ with $\KM(\gamma_0)<\varepsilon_M$, the solution of \eqref{eq:flow} satisfies
\begin{equation}\label{eq:main-convergence}
 \rho(t)^{-1}\gamma(\chi_t(\cdot),t)\longrightarrow R_\infty\beta_M
 \quad\text{in }C^\infty(\T_{\ell_M};\R^2),
\end{equation}
where $R_\infty\in SO(2)$ is a fixed planar rotation, $\chi_t:\T_{\ell_M}\to S^1$ are orientation-preserving diffeomorphisms, and
\begin{equation}\label{eq:rho-intro}
 a=L[\gamma_0]/\ell_M,\qquad \rho(t)=(a^4+30t)^{1/4}.
\end{equation}
A fixed translation $p_0\in\R^2$ can be chosen by aligning the initial curve. For every integer $j\ge0$ and every $0<\mu<g_M$, with
\begin{equation}\label{eq:certified-gaps}
 g_1=5,\qquad g_M=1/(2M^2)\quad(M\ge2),
\end{equation}
there are $\varepsilon_{M,j,\mu},C_{M,j,\mu}>0$ such that $\KM(\gamma_0)<\varepsilon_{M,j,\mu}$ implies
\begin{equation}\label{eq:main-rate}
 \left\|\frac{\gamma\circ\chi_t-p_0}{\rho(t)}-R_\infty\beta_M\right\|_{C^j}
 \le C_{M,j,\mu}\left(1+\frac{30t}{a^4}\right)^{-\mu/30}
       \sqrt{\KM(\gamma_0)},\qquad t\ge a^4.
\end{equation}
No symmetry is imposed, including between successive covering periods.
\end{theorem}

The theorem gives convergence to one fixed rotation, rather than only closeness to the rotation orbit. Its smallness hypothesis allows arbitrary smooth oscillations whose curvature-profile deviation is sufficiently small; no additional smallness in a higher Sobolev norm is assumed. The qualitative $C^\infty$ conclusion holds on one initial neighborhood, while the quantitative estimate for a prescribed $C^j$ norm and rate may require a smaller neighborhood. All constants and neighborhoods may depend on the fixed covering number $M$.

The translation $p_0$ in \eqref{eq:main-rate} is determined by the initial alignment. It is necessary even for an exact translated expander, since the profile deviation then vanishes. In the qualitative limit $p_0/\rho(t)\to0$, which explains why \eqref{eq:main-convergence} needs no translation. Exponential decay in the normalized time $\widehat t=\frac1{30}\log(1+30t/a^4)$ becomes the algebraic decay in \eqref{eq:main-rate}. In particular, the physical-time exponent reflects both the expanding scale and the stability of the normalized shape.

\subsection{The spectral result and its meaning}
Let $\mathcal A_M$ denote the normalized normal-graph linearization on $M$ periods, realized in complex $L^2$ with periodic $H^4$ domain. The second main result is that rotation generates its algebraically simple kernel and that
\begin{equation}\label{eq:gap-order}
 \gap_M:=-\sup\{\Re\lambda:\lambda\in\sigma(\mathcal A_M)\setminus\{0\}\}
 \ge g_M,\qquad \gap_M\asymp M^{-2}\quad(M\to\infty).
\end{equation}
Here $\sigma(A)$ denotes the spectrum of a closed operator $A$, $\Re\lambda$ its real part, and $a_M\asymp b_M$ means that their ratio lies between two positive constants for all sufficiently large $M$. The spectrum is discrete, and the upper gap bound has an $M$-independent constant. The lower constants $g_M$ are explicit bounds; the assertion $\gap_M\asymp M^{-2}$ specifies the sharp order in $M$, without claiming an exact leading coefficient.

The degenerating gap has a geometric interpretation. A global rotation is neutral, whereas a slowly varying rotational phase along a long covering circle is damped. Bloch decomposition turns such variations into small parameters $\alpha=j/M$. We prove the uniform bound $\Re\sigma(A_\alpha)\le-\alpha^2/2$ for $0<|\alpha|\le1/2$, and the eigenvalue issuing from rotation has the expansion $\lambda(\alpha)=-D_{\rm ph}\alpha^2+O(\alpha^4)$ with $D_{\rm ph}\ge1/2$. Evaluating this branch at $1/M$ gives the matching upper bound in \eqref{eq:gap-order}. Thus stability holds for every fixed cover, while relaxation in normalized time becomes slower as the number of traversals increases.

\subsection{Proof strategy and organization}
The normalized linearization is a fourth-order, non-self-adjoint operator. Its Jacobi elliptic coefficients have exponentially decaying Fourier expansions, which allow us to estimate the full operator by explicit row inequalities. Section~\ref{sec:geometry} derives the normalization, geometric symmetry modes and exact graph equation. Section~\ref{app:single} proves dissipation transverse to rotation and establishes the simplicity of the neutral eigenvalue.

Section~\ref{app:bloch} treats all Bloch parameters. The essential step is to remove the rotational factors exactly before approximating the adjoint vector. This keeps the error estimate uniform as $\alpha\to0$. Polynomial row weights and Bernstein positivity then control the whole parameter interval, while geometric-series estimates control every omitted Fourier mode. Appendix~\ref{app:arithmetic} specifies the scalar enclosures and finite matrices. Appendix~\ref{app:calculations} expands the geometric derivatives, nonlinear coefficients, tail reductions and finite-rank error constants. Appendix~\ref{app:expanded-data} lists all 322 finite row comparisons, the 18 rational trial coefficients, every residual coefficient and the vector-error sums. The proof assumes neither a numerically computed spectrum nor a simulation of the flow.

The nonlinear argument uses established parabolic results in their published form. In Section~\ref{sec:nonlinear}, we verify the hypotheses of the weighted local theorem of Pr\"uss, Simonett and Wilke \cite[Theorem~1.2]{PSW} to obtain positive-time smoothing from $H^2$. We then apply the normal-stability theorem of Pr\"uss, Simonett and Zacher \cite[Theorem~2.1]{PSZ} at $H^3$ regularity. Section~\ref{sec:geometric} converts curvature-profile closeness into a normal graph and returns to the original geometric flow. These abstract theorems, the global existence theorem and the standard elliptic-function identities are invoked with their hypotheses checked, rather than reproved.
